\chapter*{국문요약}
\addcontentsline{toc}{chapter}{국문요약}
\begin{center}\large\bfseries\thesistitlekr\end{center}

표면 마무리 공정에서 부품 전체의 균일 처리와 지정된 영역 주변의 국소 재작업은 서로 다른 요구를 가진다. 국소 폴리싱에서는 대상 영역에 맞추어 공구 이동, 접촉 힘, 작업 시간을 함께 구성해야 한다. 본 연구는 현재 장비로 미세 스크래치를 자율 판별하는 대신 사람이 처리 영역을 지정하는 설정을 사용한다. 공간 지시 이후에도 남는 동작·힘·시간 구성 문제에 대해, 규칙 기반 전략과 교시 재생에 비해 관측 조건부 학습이 어떤 수행 가치를 추가하는지 검토한다.

추적 기반 교시 장치로 공구 위치·자세, 3축 힘, RGB 영상을 시간적으로 대응하여 기록한다. 얼룩 중심 상대위치를 학습하고 얼룩 방향을 학습 기준에 정렬한다. Flow Matching 정책은 영상, 로봇 상태, 측정 힘 이력으로 위치·자세·목표 힘 시퀀스를 생성한다. 공통 실행 계층은 좌표 복원, 계획 연결, 명령 후처리를 수행하고 기존 하위 힘제어기가 접촉을 조절한다. 스킬 생성, 접촉 제어, 절차 완료 판정의 역할을 분리한다.

세 실험은 전략별 실행, 힘 정보 구성, 공간 재사용을 평가한다. E1에서 원 규칙 R, 힘 포함 교시 재생 T, 힘 조건부 정책 C는 각 5회에서 접촉--연속 부하 통과--이탈 절차를 완료했다. 최신 요약은 10월 19일의 추가 규칙 재실험 \(R^{*}\) 5회와 10월 15일 원 T/C를 구분하며, 원 R도 별도로 보존한다. \(R^{*}\)/T/C의 평균 완료시간은 각각 \EoneRstarTime, \EoneTTime, \EoneCTime{}초였다. T는 C보다 시간 분산과 초과 힘 적분이 작았다. \(R^{*}\)는 원 R보다 평균시간이 \RstarTimePercent\% 짧았으나, 이는 별도 추가 진단이며 실측 평균속도 정합 비교가 아니다. E2는 동작 전용 A와 힘 관측 유무에 따른 동작--힘 학습 B/C를 비교한다. A는 4/5회 완료했고, A4는 접촉 해제로 진행하지 못해 \AfourInterval{}초에 작업자가 중단했다. B/C는 각 5회 완료했으며 C의 평균 작업자 기준 완료시간은 B보다 \BCSeconds{}초(\BCPercent\%) 짧았다. 이는 E1의 자동 절차 완료시간과 다른 시간 정의다.

E3는 동일 C 체크포인트와 정규화기를 재교시·재학습 없이 다른 평면 위치·방향의 등록 작업영역에 적용한 학습 스킬의 위치·방향 전이 평가다. 저장된 전후 사진의 보드마카 잉크 면적 감소율을 공통 흑색 임계값과 보존된 가시 경계 수동 주석으로 평가했다. 본문 N/P1/D1/D2는 총 \EthreeMainN{}회이며 영상 평가 가능 시행은 \EthreeImageN{}회다. 조건별 제거율 평균\(\pm\)표본 SD는 각각 \EthreeNRemoval\(\pm\)\EthreeNSD, \EthreePoneRemoval\(\pm\)\EthreePoneSD, \EthreeDoneRemoval\(\pm\)\EthreeDoneSD, \EthreeDtwoRemoval\(\pm\)\EthreeDtwoSD\%이며 영상 유효 \(n=5,5,4,5\)다. D1-4는 D1의 5회에 남기되 최종 공구 가림으로 제거율만 산출하지 않았다. P2도 실행된 5회이며 모두 최종 가림으로 본문 주 평가에서 제외하고 원 25회 기록에 보존한다. 100\%는 해당 영상·분할 기준에서 잔여 흑색 잉크가 검출되지 않았다는 뜻이다.

결과는 인간 지정 국소 작업에서 학습 스킬의 실행 가능성, 힘 정보 구성별 실행 특성, 시험 배치 안에서 동일 정책과 작업 등록·방향 어댑터의 재사용을 뒷받침한다. 고정 작업에서는 규칙과 교시 재생도 유효했다. E3의 영상 면적 감소는 E1/E2의 제거 품질을 소급 검증하지 않으며, 금속 제거 깊이·질량·거칠기·미세 스크래치 복원은 독립 계측하지 않았다.

\vspace{0.4cm}
\noindent\textbf{주요어:} 국소 로봇 폴리싱, 인간 작업 지정, 동작--힘 스킬, 모방학습, Flow Matching, 공간 스킬 전이, 영상 기반 얼룩 제거
